% !TeX root = ../../../../../main.tex
% chapters/sections/chapter3/subsections/producers/problem_setup.tex

\subsection{生産者の問題設定}
\label{sec:chapter3_producers_problem_setup}
本モデルはヨーマン・ファーマ・モデルであるため家計\(h\)は生産者として自身の財の価格\(p_t^h\)を決定する。

\subsubsection{自国生産者の問題設定}
\label{sec:chapter3_producers_problem_setup_home}
仮定\ref{sec:chapter3_assumptions_price_optimization_home_opportunity}により
自国家計\(h\)は毎期\(1-\xi^H\)の確率で価格を改定する機会を得る。
そして価格改定の機会を得た家計\(h\)は
問題\ref{form:sec_chapter3_households_stage2_home_utility_def}において価格\(p_t^h\)を最適化する。
また\ref{itm:producer_optimization_home_calculation_remark2}より
労働\(l_t^h\)と生産\(y_t^h\)は価格\(p_t^h\)の関数であり、
\ref{itm:producer_optimization_home_calculation_assumption}より
消費者物価指数\(p_t^{H\to W}\)は価格\(p_t^h\)の関数とはみなさない。
\(t\)期で設定した価格\(p_t^h\)が\(k\)期先まで維持される確率は\((\xi^H)^k\)であることがわかり
この問題は以下の問題と同値となる（\ref{}節TODO）。
\begin{equation}
\label{form:model_producers_optimization_home_lagrangian_def}
    \begin{aligned}
        \max_{p_t^h}\operatorname{E}_t\sum_{k=0}^{\infty}(\xi^H)^k\left(\prod_{i=0}^{k-1}\beta_{t+i}^H\right)\Biggl[&\left(\log c_{t+k}^{h\to W}-\frac{\phi^H}{2}(l_{t+k}^h)^2\right) \\
        &+\lambda_{t+k}^h\biggl(\Bigl(d_{t+k}^{h\to H}+(1+i_{t+k-1}^F)\frac{e_{t+k}^{/*}}{e_{t+k-1}^{/*}}b_{t+k}^{h\to F}+(1-\tau_{t+k}^H)p_t^hy_{t+k}^h+t_{t+k}^H\Bigr) \\
        &\quad-\Bigl(\sum_{j'\in J}q_{t+k,t+k+1}(j')d_{t+k+1}^{h\to H}(j')+b_{t+k+1}^{h\to F}+p_{t+k}^{H\to W}c_{t+k}^{h\to W}\Bigr)\biggr)\Biggr]
    \end{aligned}
\end{equation}

\subsubsection{外国生産者の問題設定}
\label{sec:chapter3_producers_problem_setup_foreign}
仮定\ref{sec:chapter3_assumptions_price_optimization_foreign_opportunity}により
外国家計\(f\)は毎期\(1-\xi^F\)の確率で価格を改定する機会を得る。
そして価格改定の機会を得た家計\(f\)は
問題\ref{form:sec_chapter3_households_stage2_foreign_utility_def}において価格\(p_t^{f*}\)を最適化する。
また\ref{itm:producer_optimization_foreign_calculation_remark2}より
労働\(l_t^f\)と生産\(y_t^f\)は価格\(p_t^{f*}\)の関数であり、
\ref{itm:producer_optimization_foreign_calculation_assumption}より
消費者物価指数\(p_t^{F\to W*}\)は価格\(p_t^{f*}\)の関数とはみなさない。
\(t\)期で設定した価格\(p_t^{f*}\)が\(k\)期先まで維持される確率は\((\xi^F)^k\)であることがわかり
この問題は以下の問題と同値となる（\ref{}節TODO）。
\begin{equation}
\label{form:chapter3_producers_problem_setup_optimization_foreign_lagrangian_def}
    \begin{aligned}
        \max_{p_t^{f*}}\operatorname{E}_t\sum_{k=0}^{\infty}(\xi^F)^k\left(\prod_{i=0}^{k-1}\beta_{t+i}^F\right)\Biggl[&\left(\log c_{t+k}^{f\to W}-\frac{\phi^F}{2}(l_{t+k}^f)^2\right) \\
        &+\lambda_{t+k}^{f/*}\biggl(\Bigl(d_{t+k}^{f\to F*}+(1+i_{t+k-1}^H)\frac{e_{t+k-1}^{/*}}{e_{t+k}^{/*}}b_{t+k}^{f\to H*}+(1-\tau_{t+k}^F)p_t^{f*}y_{t+k}^f+t_{t+k}^{F*}\Bigr) \\
        &\quad-\Bigl(\sum_{j'\in J}q_{t+k,t+k+1}^*(j')d_{t+k+1}^{f\to F*}(j')+b_{t+k+1}^{f\to H*}+p_{t+k}^{F\to W*}c_{t+k}^{f\to W}\Bigr)\biggr)\Biggr]
    \end{aligned}
\end{equation}